\begin{statement}[C.8.15; Proved: every critical point and every level set]
For every zero $\mu$ of $\wp_\tau$, the zero sets of the two
critical factors are exactly
\[
 \begin{aligned}
 W=0&\quad\Longleftrightarrow\quad
       \tau\in\operatorname{SL}_2(\Z)\{i,\rho\},\\
 3W\wp'_\tau(\mu)-2g_2=0&\quad\Longleftrightarrow\quad
       \tau\in\operatorname{SL}_2(\Z)\rho.
 \end{aligned}
\]
On the coarse $t$-plane the function $L$ has no critical point
away from $t=1$. The point $t=0$ is a noncritical $C^1$ point,
with $L_t(0)=7/384$. The point $t=1$ is the unique global
maximum, of value $\log\sqrt2$, and has the fractional-power
local maximum described in C.3. In particular there is no
saddle in $0<t<1$. On the root charts, $h=0$ over $t=1$
is a nondegenerate maximum, while $\delta=0$ over $t=0$
is a degenerate critical point with nonzero harmonic cubic
leading term and three alternating ascending and descending
sectors. These are all critical points on the root charts.

For every angle $\theta$, $r\mapsto L(1+re^{i\theta})$ is
strictly decreasing for $r>0$, from $\log\sqrt2$ to $-\infty$.
For each $c<\log\sqrt2$, the level set $\{L=c\}$ is one
Jordan curve, strictly star-shaped about $1$. Its bounded
complement is $\{L>c\}$; its unbounded complement is
$\{L<c\}$. The curve is real analytic except possibly at
$t=0$ when $c=\frac14\log3$, where it is $C^1$ with nonzero
tangent. The closed superlevel sets form nested topological
discs. At the maximum the level set is the single point $1$,
and above it the level set is empty. The corresponding
statements for $D_{\rm part}$ follow by its increasing affine
change from $L$.
\end{statement}
\begin{proof}
Take a representative in the closed standard modular fundamental
region. If $T\ge7/5$, C.8.13 applies. Otherwise
$|\Re\tau|\le1/2$ and $\sqrt3/2\le T\le7/5$. For fixed
$\Re\tau$, the function
\[
 |w|^2=\frac{(\Re\tau)^2+(T-1)^2}
                    {(\Re\tau)^2+(T+1)^2}
\]
has its maximum on this interval at an endpoint and increases
with $(\Re\tau)^2$. The two endpoint bounds are $r_0^2$
and $41/601<r_0^2$. Thus C.8.14 covers the entire remaining
region. We conclude $|1-X|<1$ for every $t\ne1$.
Equations \eqref{eq:cc-radial} now prove strict radial
monotonicity, including $t=0$ by continuity. They also exclude
both critical equations $W=0$ and $3Wp-2g_2=0$ on the
ordinary base. The elliptic points must be checked on their
root charts: at $\rho$, the nonzero coarse derivative already
given rules out a critical point; at $i$, $W=p=0$ but $3Wp-2g_2=-2g_2\ne0$, and C.3 gives
the strict maximum. At $\rho$, $W=g_2=0$ and $p\ne0$,
so both displayed factors vanish. This proves their exact zero
sets, including the elliptic orbits. In particular the vanishing of both factors at
$\rho$ in the $\tau$ coordinate is a ramification effect,
not a saddle on the coarse base. More explicitly, on the root
chart at $\rho$,
\[
 L=L(0)+\frac7{192}\Re\left(\frac{ip_0^2}{\pi^3}\delta^3\right)
                   +O(|\delta|^4),
\]
with $p_0\ne0$, whereas on the chart at $i$,
$L=\log\sqrt2-\pi|h|^2/2+O(|h|^4)$.
These expansions give the asserted root-chart critical types.

At the cusp, the product formula and \eqref{eq:cc-EZ} give,
uniformly in $\Re\tau$,
\[
 L=-\frac{\pi T}{6}+\log(2\sqrt3)-
                      \frac{\pi\beta^2}{T}+O(e^{-2\pi T}),
 \qquad |t|=\frac{e^{2\pi T}}{1728}(1+o(1)).
\]
Hence $L\to-\infty$ uniformly as $|t|\to\infty$.
Continuity at $t=1$ supplies the other endpoint of each ray.
The intermediate value theorem and strict monotonicity give
one and only one radius $r_c(\theta)>0$ at each level $c$.
Properness and strict monotonicity imply continuity in $\theta$;
the nonvanishing radial derivative gives the stated local
regularity by the implicit-function theorem. Thus
$\theta\mapsto1+r_c(\theta)e^{i\theta}$ is a continuous
injective parametrization of a Jordan curve and has exactly
the asserted complementary regions.

The proposed route in the prompt that posits one saddle on
$(0,1)$ is therefore false. The corrected statement is the
strict radial foliation above. It supplies, in addition to
the separate duplication proof, a global proof of R4 and
resolves the critical-point and global-level-set question
left as D.4.6 in the supplied checkpoint.
\end{proof}
